%% ma: L11-B07-0006
%% lop: 11
%% chuong: 2
%% bai: 7
%% dang: 11.7.8
%% loai: DS
%% muc: [NB, NB, TH, VD]
%% dap_an: "ĐSĐĐ"
%% nguon: "de-on-tap-giua-ki-1-toan-11-nam-2026-2027-truong-thpt-ngo-quyen-da-nang.pdf (M09, Đề ôn tập giữa kì 1 Toán 11 THPT Ngô Quyền, Đà Nẵng), Mã đề 0111, Phần II Câu 2"
%% kien_thuc: "Bài 5 (dãy số), Bài 6 (cấp số cộng), Bài 7 (cấp số nhân)"
%% trang_thai: chinh-thuc
%% ghi_chu: "Đã duyệt 29/10/2026."
\begin{ex}
Cho dãy số $(u_n)$ có số hạng tổng quát $u_n=3n-1$ với mọi $n\in\mathbb N^*$.
\choiceTF{\True Ba số hạng đầu tiên của dãy số $(u_n)$ lần lượt là $u_1=2$, $u_2=5$, $u_3=8$}{Dãy số $(u_n)$ là một dãy số giảm}{\True Dãy số $(u_n)$ là một cấp số cộng có số hạng đầu $u_1=2$, công sai $d=3$ và tổng của $20$ số hạng đầu tiên là $S_{20}=610$}{\True Xét dãy số $(v_n)$ xác định bởi $v_n=2^{u_n}$ với mọi $n\in\mathbb N^*$. Khi đó $(v_n)$ là một cấp số nhân có số hạng đầu $v_1=4$ và công bội $q=8$}
\loigiai{a) Đúng.

b) Sai: $u_{n+1}-u_n=3>0$ nên dãy tăng.

c) Đúng: $d=3$, $S_{20}=\dfrac{20(u_1+u_{20})}2=10(2+59)=610$.

d) Đúng: $v_1=2^2=4$, $\dfrac{v_{n+1}}{v_n}=2^{u_{n+1}-u_n}=2^3=8$.}
\end{ex}
